\section{A radial estimate for the larger discriminants}\label{app:radial}
For the range $D\geq875$ we use an entrywise estimate in place of the
spectral estimate of Section~\ref{sec:metric}. The two estimates start from
the same normalized matrix. Here the additional saving comes from keeping
track of which nodes and derivative orders can occur together in a
determinant term. A node used many times must supply smaller derivative
orders, while nodes near the origin give stronger radial decay. We express
these restrictions by a convex function of the numbers of columns selected
at the different radii.

Throughout this appendix assume $K\geq2$ and $S\geq1$. Put $R=KS$, $M=R(S+1)/2$, $W=R^2S/3$, and
$\theta=M/(nJ)\leq1$. The derivative scale is denoted by $\ell$, as in
\eqref{eq:residual-matrix}. We assume $\ell\geq0$ and put
$\varepsilon=J/R$. The parameters in Appendix~\ref{app:global} satisfy
these assumptions.

\subsection{Counting the ordered exponents}
Put
\[
 u(G)=1-\sqrt{1-G},\qquad
 v(G)=\begin{cases}
 \sqrt{G/K},&0\leq G\leq1/K,\\
 1-\sqrt{(1-1/K)(1-G)},&1/K\leq G\leq1.
 \end{cases}
\]
The second function is the quantile of $r+s$, after scaling by $R$, on
the continuous triangle $r+Ks\leq R$. The following bounds account for
the discrete rows of the interpolation matrix.

\begin{lemma}\label{lem:radial-quantiles}
List the values of $r$ and $r+s$ on $\Lambda$ in nondecreasing order as
$r_{(k)}$ and $z_{(k)}$, respectively. For $G_k=k/M$,
\[
 Ru(G_k)-1\leq r_{(k)}\leq(R+K)u(G_k),\qquad
 z_{(k)}\geq Rv(G_k)-2.
\]
\end{lemma}
\begin{proof}
For $0\leq a\leq R$, the number of rows with $r\geq a$, with $a$ an
integer, is
\[
 T_r(a)=\sum_{b=1}^S(Kb-a)_+
       =\frac K2\{y^2+y+\alpha(1-\alpha)\},
 \quad y=S-a/K,\quad\alpha=\{a/K\}.
\]
This follows by summing the arithmetic progression of positive terms.
It implies
\[
 M(1-a/R)^2\leq T_r(a)
       \leq M(1-a/(R+K))^2.
\]
For the first inequality, subtracting its right-hand quadratic from
$y^2+y$ leaves $y(1-y/S)\geq0$.
For the second, write $a/K=q+\alpha$. After multiplication by $S+1$,
the difference of the two sides, apart from the positive factor $K/2$,
is $q(S+1-q-2\alpha)+S\alpha^2\geq0$.
The endpoints $a=0,R$ satisfy the same inequalities.
Applying the lower tail bound at $a=r_{(k)}+1$ and the upper tail bound
at $a=r_{(k)}$ proves the two assertions about $r_{(k)}$.

For the joint exponent, let $N(t)$ count rows with $r+s\leq t$, for an
integer $0\leq t<R$. Let $F$ be the inverse of $v$, extended by $F(x)=1$
for $x\geq1$. We show that $N(t)\leq MF((t+2)/R)$.
If $t\leq S-2$, the whole triangle $r+s\leq t$ is contained in
$\Lambda$, so
\[
 N(t)=\frac{(t+1)(t+2)}2
       \leq\frac{S+1}{2S}(t+2)^2=MF((t+2)/R).
\]
If $t\geq S-1$, put $h=K-1$ and $A=R-t-1$.
For $A\geq1$ the complementary count is
\[
 M-N(t)=\sum_{s=0}^{S-1}(A-hs)_+
       \geq\frac{A^2}{2h}+\frac A2.
\]
Here $1\leq A\leq hS$. On this interval
\[
 A^2+hA\geq(1+1/S)(A-1)^2:
\]
the difference is a concave quadratic in $A$ and is nonnegative at both
endpoints. The last lower bound is therefore at least
$M\{1-F((t+2)/R)\}$. When $A=0$ the desired assertion is immediate.
Taking $t=z_{(k)}$ and using $k\leq N(t)$ proves the remaining bound.
\end{proof}

\subsection{The falling factorial and the derivative orders}
For $0\leq t\leq x\leq1$, define
\[
 f(x,t)=x\log x-(x-t)\log(x-t)-t,
 \qquad
 g(x,t)=\max_{0\leq a\leq\min(x,t)}\{a\ell+f(x,a)\},
\]
with $0\log0=0$. We also use $g(x,t)$ for $t>x$ with this same
definition. Its maximum occurs at
$a=\max(0,\min(t,x-e^{-\ell}))$.
For $0\leq j\leq r<R$, integral comparison gives
\begin{equation}\label{eq:radial-factorial}
 \log\frac{\fall rj}{R^j}
 \leq Rf\left(\frac{r+1}{R},\frac jR\right).
\end{equation}
Indeed, on each unit interval the decreasing function
$\log((r+1-t)/R)$ is at least its value at the right endpoint.
This also proves the inequality when $j=r$.

Fix a set $C$ of $M$ columns, and write $c_i$ for the number selected at
node $i$. Replace its selected orders by $J-c_i,\ldots,J-1$, and list
the resulting multiset increasingly as $j_k^+$. If the original orders
are listed as $j_{(k)}$, then
\begin{equation}\label{eq:radial-order-envelope}
 j_{(k)}\leq j_k^+\leq t_k,
 \qquad t_k=J-\frac{M-k}{n}.
\end{equation}
The first assertion follows by ordering the columns at each node; the
second follows because at most $n(J-a)$ orders can be at least $a$.

Consider a nonzero term in the expansion of a selected determinant.
Every assigned derivative satisfies $j\leq r$. The mixed derivative
$f_{xt}=1/(x-t)$ is positive. Exchanging two oppositely ordered pairs
therefore increases the sum of $f$; the exchange preserves feasibility,
since the larger derivative was already assigned to the smaller row degree.
Successive exchanges pair the increasingly ordered row degrees and
derivatives. Increasing the row bound and then using that $g$ is
nondecreasing in its second argument gives
\begin{equation}\label{eq:radial-rearrangement}
 \sum_C\left(j\ell+
       \log\frac{\fall rj}{R^j}\right)
 \leq R\sum_{k=1}^M g(x_k,j_k^+/R),
 \qquad x_k=\min\left(1,(1+1/S)u(G_k)+1/R\right).
\end{equation}
The exchange argument extends to the boundary $x=t$ by continuity.

Put $h(x,t)=t\ell-g(x,t)$. For $\ell\geq0$, this function is
nonnegative, increases with $t$ and $\ell$, and decreases with $x$.
For example $g_t=(\ell+\log(x-t))_+$ before the maximum becomes
constant, so $0\leq g_t\leq\ell$. The other assertions follow either
by differentiation at the maximizing argument or directly from the
definition. Thus an admissible normalized factorial saving is any
number $S_{\rm der}$ satisfying
\begin{equation}\label{eq:radial-derivative-saving}
 0\leq S_{\rm der}\leq\frac3M\sum_{k=1}^M h(x_k,t_k/R).
\end{equation}
It can be bounded by grouping consecutive indices: on a group with
$x_k\leq b$ and $t_k/R\geq a$, use its relative size times $3h(b,a)$.
This requires only a finite sum, rather than a list of the $M$ rows.

\subsection{An additional saving from the number of columns at a node}
Suppose $\theta\varepsilon<1/2$. Choose integers
$0<m_1<\cdots<m_L<M/n$, put $\tau_r=m_r/J$, and suppose
\[
 a_r=1-\sqrt{\tau_r/\theta}-(1-\tau_r)\varepsilon>0,
 \qquad \nu_r=\ell+\log a_r>0.
\]
These $\nu_r$ decrease with $r$: the derivative of
$1-\sqrt{\tau/\theta}-(1-\tau)\varepsilon$ is
$\varepsilon-1/(2\sqrt{\theta\tau})<0$ for $\tau<\theta$.
Set $\nu_{L+1}=0$ and $\lambda_r=\nu_r-\nu_{r+1}$.
Smaller nonnegative increments may be used throughout.

If $t_k\geq J-m_r$, then $G_k\geq1-\tau_r/\theta$.
The function
$u(G)-\varepsilon(1-\theta(1-G))$ is increasing, since its derivative
is at least $1/2-\theta\varepsilon>0$.
It follows that $x_k-t_k/R\geq a_r$. On the interval between
$j_k^+$ and $t_k$, the derivative of $Rg(x_k,t/R)$ is therefore at
least the staircase
$\sum_r\lambda_r\boldsymbol1_{t\geq J-m_r}$.
Integration gives an improvement to \eqref{eq:radial-rearrangement}.
The improvement can be summed explicitly, because
\begin{align*}
 \sum_k(t_k-J+m)_+&=\frac{nm^2+m}{2},\\
 \sum_k(j_k^+-J+m)_+
 &=\frac{nm(m-1)}2
       -\frac12\sum_i(m-c_i)_+(m-c_i-1).
\end{align*}
Subtracting these identities, and using
$\sum_i(m-c_i)_+\leq nm$, gives a difference at least
$\frac12\sum_i(m-c_i)_+^2$. Consequently
\begin{equation}\label{eq:radial-derivative-coupled}
 \frac2W\sum_C\left(j\ell+\log\frac{\fall rj}{R^j}\right)
 \leq\gamma\ell-S_{\rm der}
       -\frac1W\sum_{i,r}\lambda_r(m_r-c_i)_+^2,
 \quad
 \gamma=\frac2W\left(MJ(1-\theta/2)+\theta J/2\right).
\end{equation}
In obtaining the factorial saving we used $2RM/W=3(1+1/S)\geq3$.

\subsection{Radial weights and the convex estimate}
Order the nodes by modulus. Suppose $N$ divides $n$, and divide the
nodes into $N$ consecutive groups of $n/N$ nodes. Choose numbers
$L_0\geq\cdots\geq L_{N-1}\geq L_N=0$ such that
\[
 \log(\rho/|z_i|)\geq L_b\qquad
       (bn/N<i\leq(b+1)n/N).
\]
At a zero node the assertion has its limiting meaning. The disk estimate
\eqref{eq:node-disk} permits the choice
\begin{equation}\label{eq:radial-bin-weights}
 L_b=\log\frac{1+\eta}
 {\sqrt{\min(1,(b+1)/N+1/n)}+\eta},
 \qquad \eta=\frac{\sqrt{D+4}/4}{\sqrt{n\sqrt D/(2\pi)}}.
\end{equation}
The extra $1/n$ weakens the bound and is convenient when a common bound
is used for several values of $n$.
Choose an index $b_0$ and a number $\overline L\geq L_{b_0}$.
Write
\[
 w_b=\frac{N}{nJ}\sum_{i\text{ in group }b}c_i,
 \qquad G_b(w)=\frac1{N\theta}\sum_{a=0}^bw_a.
\]
Thus $0\leq w_b\leq1$ and $\sum_bw_b=N\theta$.
Define
\begin{align}
 \Phi(G)&=3\int_0^G\left(v(t)-\frac{J+2}{R}\right)_+\,dt,
                                                        \label{eq:radial-Phi}\\
 \Delta&=\frac2W\left\{\binom{J+2}{3}
                         -\binom{(J-S)_++2}{3}\right\}.
                                                        \label{eq:radial-Delta}
\end{align}
In particular, when $J\geq S$,
\[
 \Delta=3\left(\frac{J+1}{R}\right)^2
       -\frac3K\frac{J+1}{R}+\frac1{K^2}-\frac1{R^2}.
\]
On the polytope $\mathcal P_\theta=\{w\in[0,1]^N:\sum_bw_b=N\theta\}$,
put
\begin{align}
 \mathcal F(w)={}&\sum_{b=0}^{N-1}(L_b-L_{b+1})\Phi(G_b(w))
                                                     \notag\\
 &+\frac{3\varepsilon}{2N\theta}\sum_{b=0}^{N-1}
   \left\{L_{\max(b,b_0)}w_b^2
             +\sum_{r=1}^L\lambda_r(\tau_r-w_b)_+^2\right\}
       -\overline L\Delta.                 \label{eq:radial-functional}
\end{align}

\begin{proposition}[Radial determinant estimate]\label{prop:radial-size}
With the preceding choices, let $F_*\leq\min_{\mathcal P_\theta}\mathcal F$.
Then every maximal minor satisfies
\begin{equation}\label{eq:radial-minor-bound}
 \frac1W\log|\det B_C(\ell)|^2
 \leq\gamma\ell-S_{\rm der}-F_*+\frac{2M\log M}{W}.
\end{equation}
\end{proposition}
\begin{proof}
Consider a nonzero determinant term, and put $z=r+s$ for its assigned
row. Since $z\geq r\geq j$, for a radial weight $l\geq0$ and
$\widetilde l=\min(l,\overline L)$ one has
\[
 (z-j)l\geq(z-J)_+l+(J-j)\widetilde l
                               -(J-z)_+\widetilde l.
\]
Indeed, subtracting the first term leaves
$(\min(z,J)-j)l$, whose coefficient is nonnegative.
This observation allows us to retain the full radial weight on the
first term while bounding the possible negative correction by
$\overline L$.

Among all assignments of rows to columns, pairing the smallest values
of $(r+s-J)_+$ with the largest radial weights minimizes their sum.
For any integer $h\leq M$, Lemma~\ref{lem:radial-quantiles} and the
right-endpoint integral inequality for an increasing function give
\[
 \frac2W\sum_{k=1}^h(z_{(k)}-J)_+
 \geq\frac{2RM}{3W}\Phi(h/M)\geq\Phi(h/M).
\]
Summation by parts over the node groups therefore gives the first sum
in \eqref{eq:radial-functional}.
At node $i$, the distinct derivative orders satisfy
$\sum(J-j)\geq c_i(c_i+1)/2\geq c_i^2/2$.
Also
\[
 \sum_{\Lambda}(J-r-s)_+
 \leq\sum_{s=0}^{S-1}\binom{(J-s)_++1}{2}
 =\binom{J+2}{3}-\binom{(J-S)_++2}{3}.
\]
Thus the clipped radial terms contribute at least
$W^{-1}\sum_i c_i^2L_{\max(b(i),b_0)}-\overline L\Delta$.

Combine this with \eqref{eq:radial-derivative-coupled}.
Convexity on each node group gives
\[
 \sum_{i\text{ in group }b}c_i^2\geq\frac nN J^2w_b^2,
 \qquad
 \sum_{i\text{ in group }b}(m_r-c_i)_+^2
       \geq\frac nN J^2(\tau_r-w_b)_+^2.
\]
Since $nJ^2/W=3\varepsilon(1+1/S)/(2\theta)$, dropping the factor
$1+1/S$ proves that the total normalized saving is at least
$S_{\rm der}+\mathcal F(w)$. The modulus of each nonzero determinant
term is bounded accordingly. There are at most $M!\leq M^M$ terms;
taking the square of their sum proves \eqref{eq:radial-minor-bound}.
\end{proof}

The function $\Phi$ is convex, because its derivative is increasing.
Every coefficient of the convex terms in \eqref{eq:radial-functional}
is nonnegative. Hence $\mathcal F$ is convex. In particular, any
$w^*\in\mathcal P_\theta$ gives the explicit lower bound
\begin{equation}\label{eq:radial-supporting-plane}
 F_*=\mathcal F(w^*)-\nabla\mathcal F(w^*)\cdot w^*
                 +\min_{w\in\mathcal P_\theta}
                         \nabla\mathcal F(w^*)\cdot w.
\end{equation}
The last minimum is found by filling the coordinates with the smallest
coefficients first, subject to the upper bound one and the prescribed
sum. This follows by transferring mass from a larger coefficient to a
smaller one whenever the latter coordinate is not full. No assertion
that $w^*$ minimizes $\mathcal F$ is required.
For numerical evaluation its gradient is obtained directly from
\eqref{eq:radial-functional}; the positive-part squares are continuously
differentiable, including at their endpoints.
